\section{Separating Vision from Language}
\label{sec:vision-language}

Question-level accuracy alone cannot tell us whether a model's spatial answers are genuinely grounded in the image or are largely recoverable from object names, task format, and generic priors. Ordinary-Bench therefore includes a second layer of analysis that explicitly separates visual contribution from language-driven guessing. The benchmark is designed so that the same scene, the same probe set, and the same scoring pipeline can be evaluated under controlled perturbations of the visual input.

\subsection{Three-Condition Protocol}
\label{sec:three-condition}

Section~\ref{sec:protocol} defined the three evaluation conditions: correct image, mismatched image, and no image. Their role is not merely ablative, but diagnostic. Condition~A asks whether the model can use the appropriate visual evidence. Condition~B tests whether the model is sensitive to image-content mismatch rather than simply exploiting the presence of \emph{some} image. Condition~C estimates how much of the task can be solved from textual priors alone.

This setup yields a scene-level measure of visual grounding. Let $\mathrm{NRMS}_A$, $\mathrm{NRMS}_B$, and $\mathrm{NRMS}_C$ denote reconstruction error under the three conditions. We define
\begin{equation}
  \mathrm{VIG} = \mathrm{NRMS}_C - \mathrm{NRMS}_A,
\end{equation}
so that a positive score indicates that the correct image leads to a better recovered scene than text-only prompting. We also track the gap between mismatched and correct images, $\mathrm{NRMS}_B - \mathrm{NRMS}_A$, to test whether the model is using scene-specific visual evidence rather than merely benefiting from multimodal formatting. In the final evaluation, these quantities should be paired with scene-wise significance tests such as permutation tests or paired nonparametric tests.

\subsection{Consistency as Structural Evidence}
\label{sec:consistency}

Visual grounding should improve not only accuracy but also the internal structure of a model's answers. We therefore complement VIG with consistency diagnostics that operate directly on the predicted constraint set.

\paragraph{Transitivity for QRR.}
Distance comparisons admit natural transitive checks. If a model predicts $d(A,B) < d(C,D)$ and $d(C,D) < d(E,F)$, then a coherent ordinal belief should also support $d(A,B) < d(E,F)$. We summarize violations with a \emph{transitivity violation rate},
\begin{equation}
  T_{\mathrm{viol}} = \frac{\#\{\text{violated testable transitive implications}\}}{\#\{\text{testable transitive implications}\}}.
\end{equation}
Low transitivity violation indicates that the model's distance judgments are compatible with a single partial order rather than a collection of unrelated local guesses.

\paragraph{Reciprocity for TRR.}
Directional judgments admit an analogous reciprocity check. If a target is judged to lie at 3~o'clock in the reference frame defined by $(\texttt{ref1}, \texttt{ref2})$, then swapping the viewing direction should induce the corresponding opposite direction. We summarize such failures with a \emph{reciprocity violation rate} $R_{\mathrm{viol}}$. These structural metrics are especially informative because purely language-driven heuristics may occasionally answer individual probes correctly while still failing to satisfy the geometric relations that tie probes together.

\paragraph{Selective answering.}
The benchmark is also compatible with selective-answering variants in which the model may abstain, return a coarse answer, or express ambiguity when the visual evidence is weak. In that setting, abstention is not treated as failure by default; rather, it can improve the quality of the extracted constraint set by preventing unreliable observations from contaminating reconstruction. This follows the broader perspective that reliable multimodal evaluation should distinguish ``cannot tell'' from ``guessed incorrectly'' \cite{whitehead2022reliablevqa,eisenschlos2024selectivevqa,khan2024consistency}.

\subsection{Random Baselines and Belief Faithfulness}
\label{sec:baselines}

To calibrate both accuracy and reconstruction quality, we compare model behavior against random baselines matched to task arity: $1/3$ for QRR, $1/12$ for TRR hour-level accuracy, and $1/4$ for TRR quadrant accuracy. However, random-answer baselines are not sufficient on their own, because the reconstruction solver may impose weak geometric regularities even on noisy constraint sets. We therefore recommend reporting a \emph{belief faithfulness gap}: the difference between the model's reconstruction quality and that of a random baseline under the same solver, metrics, and scene set.

This comparison matters for interpretation. A model with only modest question-level gains over random may still produce a substantially more coherent global belief, while a model with apparently non-trivial accuracy may collapse to near-random behavior once its answers are examined structurally. Ordinary-Bench is intended to make this distinction explicit.
